% Options for packages loaded elsewhere
% Options for packages loaded elsewhere
\PassOptionsToPackage{unicode}{hyperref}
\PassOptionsToPackage{hyphens}{url}
\PassOptionsToPackage{dvipsnames,svgnames,x11names}{xcolor}
%
\documentclass[
  spanish,
  13pt,
  letterpaper,
  DIV=11,
  numbers=noendperiod]{scrartcl}
\usepackage{xcolor}
\usepackage[margin=2.5cm]{geometry}
\usepackage{amsmath,amssymb}
\setcounter{secnumdepth}{5}
\usepackage{iftex}
\ifPDFTeX
  \usepackage[T1]{fontenc}
  \usepackage[utf8]{inputenc}
  \usepackage{textcomp} % provide euro and other symbols
\else % if luatex or xetex
  \usepackage{unicode-math} % this also loads fontspec
  \defaultfontfeatures{Scale=MatchLowercase}
  \defaultfontfeatures[\rmfamily]{Ligatures=TeX,Scale=1}
\fi
\usepackage{lmodern}
\ifPDFTeX\else
  % xetex/luatex font selection
\fi
% Use upquote if available, for straight quotes in verbatim environments
\IfFileExists{upquote.sty}{\usepackage{upquote}}{}
\IfFileExists{microtype.sty}{% use microtype if available
  \usepackage[]{microtype}
  \UseMicrotypeSet[protrusion]{basicmath} % disable protrusion for tt fonts
}{}
\makeatletter
\@ifundefined{KOMAClassName}{% if non-KOMA class
  \IfFileExists{parskip.sty}{%
    \usepackage{parskip}
  }{% else
    \setlength{\parindent}{0pt}
    \setlength{\parskip}{6pt plus 2pt minus 1pt}}
}{% if KOMA class
  \KOMAoptions{parskip=half}}
\makeatother
% Make \paragraph and \subparagraph free-standing
\makeatletter
\ifx\paragraph\undefined\else
  \let\oldparagraph\paragraph
  \renewcommand{\paragraph}{
    \@ifstar
      \xxxParagraphStar
      \xxxParagraphNoStar
  }
  \newcommand{\xxxParagraphStar}[1]{\oldparagraph*{#1}\mbox{}}
  \newcommand{\xxxParagraphNoStar}[1]{\oldparagraph{#1}\mbox{}}
\fi
\ifx\subparagraph\undefined\else
  \let\oldsubparagraph\subparagraph
  \renewcommand{\subparagraph}{
    \@ifstar
      \xxxSubParagraphStar
      \xxxSubParagraphNoStar
  }
  \newcommand{\xxxSubParagraphStar}[1]{\oldsubparagraph*{#1}\mbox{}}
  \newcommand{\xxxSubParagraphNoStar}[1]{\oldsubparagraph{#1}\mbox{}}
\fi
\makeatother


\usepackage{longtable,booktabs,array}
\newcounter{none} % for unnumbered tables
\usepackage{calc} % for calculating minipage widths
% Correct order of tables after \paragraph or \subparagraph
\usepackage{etoolbox}
\makeatletter
\patchcmd\longtable{\par}{\if@noskipsec\mbox{}\fi\par}{}{}
\makeatother
% Allow footnotes in longtable head/foot
\IfFileExists{footnotehyper.sty}{\usepackage{footnotehyper}}{\usepackage{footnote}}
\makesavenoteenv{longtable}
\usepackage{graphicx}
\makeatletter
\newsavebox\pandoc@box
\newcommand*\pandocbounded[1]{% scales image to fit in text height/width
  \sbox\pandoc@box{#1}%
  \Gscale@div\@tempa{\textheight}{\dimexpr\ht\pandoc@box+\dp\pandoc@box\relax}%
  \Gscale@div\@tempb{\linewidth}{\wd\pandoc@box}%
  \ifdim\@tempb\p@<\@tempa\p@\let\@tempa\@tempb\fi% select the smaller of both
  \ifdim\@tempa\p@<\p@\scalebox{\@tempa}{\usebox\pandoc@box}%
  \else\usebox{\pandoc@box}%
  \fi%
}
% Set default figure placement to htbp
\def\fps@figure{htbp}
\makeatother



\ifLuaTeX
\usepackage[bidi=basic,shorthands=off]{babel}
\else
\usepackage[bidi=default,shorthands=off]{babel}
\fi
\ifLuaTeX
  \usepackage{selnolig} % disable illegal ligatures
\fi


\setlength{\emergencystretch}{3em} % prevent overfull lines

\providecommand{\tightlist}{%
  \setlength{\itemsep}{0pt}\setlength{\parskip}{0pt}}



 


\KOMAoption{captions}{tableheading}
\makeatletter
\@ifpackageloaded{caption}{}{\usepackage{caption}}
\AtBeginDocument{%
\ifdefined\contentsname
  \renewcommand*\contentsname{Tabla de contenidos}
\else
  \newcommand\contentsname{Tabla de contenidos}
\fi
\ifdefined\listfigurename
  \renewcommand*\listfigurename{Listado de Figuras}
\else
  \newcommand\listfigurename{Listado de Figuras}
\fi
\ifdefined\listtablename
  \renewcommand*\listtablename{Listado de Tablas}
\else
  \newcommand\listtablename{Listado de Tablas}
\fi
\ifdefined\figurename
  \renewcommand*\figurename{Figura}
\else
  \newcommand\figurename{Figura}
\fi
\ifdefined\tablename
  \renewcommand*\tablename{Tabla}
\else
  \newcommand\tablename{Tabla}
\fi
}
\@ifpackageloaded{float}{}{\usepackage{float}}
\floatstyle{ruled}
\@ifundefined{c@chapter}{\newfloat{codelisting}{h}{lop}}{\newfloat{codelisting}{h}{lop}[chapter]}
\floatname{codelisting}{Listado}
\newcommand*\listoflistings{\listof{codelisting}{Listado de Listados}}
\makeatother
\makeatletter
\makeatother
\makeatletter
\@ifpackageloaded{caption}{}{\usepackage{caption}}
\@ifpackageloaded{subcaption}{}{\usepackage{subcaption}}
\makeatother
\usepackage{bookmark}
\IfFileExists{xurl.sty}{\usepackage{xurl}}{} % add URL line breaks if available
\urlstyle{same}
\hypersetup{
  pdftitle={Reporte de Reproducibilidad},
  pdfauthor={Autor 1; Autor 2; Autor 3},
  pdflang={es},
  colorlinks=true,
  linkcolor={blue},
  filecolor={Maroon},
  citecolor={Blue},
  urlcolor={Blue},
  pdfcreator={LaTeX via pandoc}}


\title{Reporte de Reproducibilidad}
\usepackage{etoolbox}
\makeatletter
\providecommand{\subtitle}[1]{% add subtitle to \maketitle
  \apptocmd{\@title}{\par {\large #1 \par}}{}{}
}
\makeatother
\subtitle{Artículo: {[}título del artículo a reproducir{]}}
\author{Autor 1 \and Autor 2 \and Autor 3}
\date{2026-04-23}
\begin{document}
\maketitle

\renewcommand*\contentsname{Tabla de contenidos}
{
\setcounter{tocdepth}{3}
\tableofcontents
}

\section{Selección de
artículo/reporte}\label{selecciuxf3n-de-artuxedculoreporte}

\begin{itemize}
\tightlist
\item
  Referencia del artículo seleccionado
\item
  Breve resumen del artículo seleccionado
\item
  Justificación de la selección del artículo
\end{itemize}

\section{Evaluación de
reproducibilidad}\label{evaluaciuxf3n-de-reproducibilidad}

Análisis detallado del artículo seleccionado, evaluando los siguientes
aspectos:

\begin{itemize}
\tightlist
\item
  \textbf{Disponibilidad de datos}: ¿El artículo proporciona acceso a
  los datos utilizados en el estudio? Si no es así, ¿qué dificultades se
  enfrentaron para obtener los datos?
\item
  \textbf{Disponibilidad de código}: ¿El artículo proporciona acceso al
  código utilizado para el análisis? Si es así, ¿cómo se puede acceder?
\item
  \textbf{Documentación}: ¿El artículo incluye una descripción clara de
  los métodos y procedimientos utilizados en el estudio? ¿Es suficiente
  para que otro investigador pueda replicar el estudio?
\item
  \textbf{Transparencia}: ¿El artículo proporciona información sobre los
  posibles conflictos de interés, financiamiento, y otros aspectos
  relacionados con la transparencia de la investigación?
\end{itemize}

\section{Análisis reproducible}\label{anuxe1lisis-reproducible}

Seleccionar un resultado específico del artículo (por ejemplo, una tabla
de resultados o una figura) y tratar de reproducirlo utilizando los
datos y el código proporcionados. Documentar el proceso de reproducción,
incluyendo cualquier dificultad encontrada y cómo se resolvió.

\subsection{Resultado a reproducir}\label{resultado-a-reproducir}

\begin{itemize}
\tightlist
\item
  Imagen de tabla(s), figura(s) o resultado específico del artículo
  seleccionado, junto a una breve descripción.
\end{itemize}

\subsection{Proceso de reproducción}\label{proceso-de-reproducciuxf3n}

\begin{itemize}
\tightlist
\item
  Descripción detallada del proceso seguido para reproducir el resultado
  seleccionado, incluyendo los pasos realizados, las herramientas
  utilizadas, y cualquier ajuste necesario para lograr la reproducción.
  Incluye las siguientes sub secciones:
\end{itemize}

\subsubsection{Procesamiento}\label{procesamiento}

\begin{itemize}
\item
  Incluye el código utilizado para procesar los datos, junto con una
  explicación de cada paso.
\item
  Los datos originales deben ser llamados desde input/data/original, el
  código de procesamiento se debe incluir en este mismo archivo
  (reporte-repro.qmd).
\end{itemize}

\subsubsection{Reproducción}\label{reproducciuxf3n}

\begin{itemize}
\tightlist
\item
  Generación de tabla o figura a reproducir.
\end{itemize}

\section{Conclusiones}\label{conclusiones}

Basándose en el análisis realizado, escriban una conclusión sobre el
nivel de reproducibilidad del artículo seleccionado.

\section{Recomendaciones}\label{recomendaciones}

Recomendaciones para mejorar la reproducibilidad del artículo, si es
necesario.

\section{Referencias}\label{referencias}

\protect\phantomsection\label{refs}

\section{Apéndice}\label{apuxe9ndice}

\subsection{Material suplementario}\label{material-suplementario}

\begin{itemize}
\tightlist
\item
  Incluyan aquí detalles metodológicos adicionales, en caso de ser
  necesario
\end{itemize}

\subsection{Código}\label{cuxf3digo}

\begin{itemize}
\tightlist
\item
  Incluir acá el código original, en caso de estar disponible.
\end{itemize}




\end{document}
